% ==============================================================================
% TCC - Nome do Aluno
% Capítulo 2 - Referencial Teórico
% ==============================================================================
\chapter{Fundamentação Teórica e Tecnologias Utilizadas}
\label{sec-fundteo}

Este capítulo apresenta os fundamentos teóricos e as tecnologias que
sustentam o desenvolvimento deste trabalho. A Seção~\ref{sec-fundteo-rag}
introduz o conceito de \textit{Retrieval-Augmented Generation} (RAG) e discute
as limitações da abordagem vetorial tradicional. A
Seção~\ref{sec-fundteo-graphrag} apresenta o RAG baseado em grafos de
conhecimento, com ênfase na arquitetura, no \textit{pipeline} de indexação e
nas modalidades de consulta implementadas pela biblioteca GraphRAG,
desenvolvida pela Microsoft, que será utilizada como núcleo deste trabalho. A
Seção~\ref{sec-fundteo-baseline} discute a abordagem de RAG vetorial que será
utilizada como \textit{baseline} de comparação. Por fim, a
Seção~\ref{sec-fundteo-stack} apresenta as tecnologias auxiliares escolhidas
para o desenvolvimento da aplicação web que compõe a prova de conceito.

%%% Início de seção. %%%
\section{\textit{Retrieval-Augmented Generation} (RAG)}
\label{sec-fundteo-rag}

Modelos de linguagem de grande porte (\textit{Large Language Models} -- LLMs) \cite{vaswani2017attention, brown2020language}
são treinados sobre grandes volumes de texto até uma determinada data de
corte, o que os torna limitados quando confrontados com perguntas que exigem
conhecimento posterior a esse treinamento ou conhecimento específico de um
domínio não representado de forma significativa nos dados de treinamento. Além disso, LLMs estão
sujeitos ao fenômeno de \textit{alucinação}, no qual o modelo gera respostas
plausíveis, porém factualmente incorretas ou não fundamentadas em nenhuma
fonte real.

O paradigma de \textit{Retrieval-Augmented Generation} \cite{lewis2020retrieval}
foi proposto como forma de mitigar essas limitações, combinando um mecanismo
de recuperação de informação (\textit{retrieval}) a um modelo de geração de
texto. Em vez de depender exclusivamente do conhecimento paramétrico
codificado nos pesos do modelo durante o treinamento, um sistema RAG recupera,
no momento da consulta, os trechos mais relevantes de uma base de
conhecimento externa e os inclui no contexto (\textit{prompt}) fornecido ao
LLM, que então gera a resposta com base nesse conteúdo recuperado. Esse
mecanismo permite atualizar o conhecimento do sistema sem a necessidade de
re-treinar o modelo, além de tornar as respostas passíveis de rastreamento até
sua fonte.

Na implementação mais comum de RAG, conhecida como RAG vetorial ou RAG
ingênuo (\textit{naive RAG}), o processo ocorre em duas etapas principais. Na
etapa de indexação, os documentos da base de conhecimento são divididos em
fragmentos menores (\textit{chunks}), e cada fragmento é convertido em um
vetor numérico de alta dimensão por meio de um modelo de
\textit{embeddings}, que captura sua representação semântica. Esses vetores
são armazenados em um banco de dados vetorial \cite{johnson2019billion}. Na etapa de consulta, a
pergunta do usuário é igualmente convertida em vetor, e os fragmentos
armazenados com maior similaridade de \textit{cosseno} em relação a esse vetor
são recuperados e inseridos no contexto do LLM.

Embora essa abordagem seja eficaz em diversos cenários, ela apresenta
limitações relevantes para o presente trabalho. A recuperação por similaridade
vetorial opera sobre fragmentos de texto de forma isolada, sem representar
explicitamente as relações estruturais entre diferentes partes do corpus. Em
domínios como o jurídico \cite{cui2023chatlaw}, nos quais um artigo de lei frequentemente depende de
definições, exceções e remissões estabelecidas em outros artigos, essa
limitação pode levar à recuperação de fragmentos parcialmente relevantes, mas
insuficientes para compor uma resposta completa. Adicionalmente, o RAG
vetorial tradicional apresenta dificuldade em responder perguntas de natureza
global ou holística -- por exemplo, ``quais são os principais direitos do
consumidor previstos no código?'' --, uma vez que a resposta demandaria a
síntese de informações dispersas por praticamente todo o corpus, e não a
recuperação de um pequeno conjunto de fragmentos localizados.

%%% Início de seção. %%%
\section{RAG Baseado em Grafos de Conhecimento (GraphRAG)}
\label{sec-fundteo-graphrag}

Como alternativa às limitações do RAG vetorial, diferentes abordagens têm
proposto a estruturação da base de conhecimento como um grafo, no qual nós
representam entidades do domínio e arestas representam as relações entre
elas. Essa representação preserva explicitamente a estrutura relacional do
corpus e permite estratégias de recuperação capazes de percorrer múltiplos
saltos (\textit{multi-hop}) entre entidades relacionadas, além de viabilizar a
geração de resumos hierárquicos sobre agrupamentos de entidades (comunidades),
o que favorece a resposta a perguntas de natureza global.

\subsection{O \textit{framework} GraphRAG da Microsoft}
\label{sec-fundteo-graphrag-ms}

Este trabalho utiliza o GraphRAG, biblioteca de código aberto desenvolvida
pela Microsoft Research \cite{edge2024global}, disponibilizada como pacote Python (\texttt{graphrag}) e distribuída sob
licença MIT. O GraphRAG propõe um método estruturado e hierárquico de RAG
que, diferentemente da abordagem vetorial ingênua, extrai um grafo de
conhecimento a partir do texto bruto por meio de um LLM, constrói uma
hierarquia de comunidades sobre esse grafo e gera resumos textuais para cada
comunidade, utilizando essas estruturas -- e não apenas fragmentos de texto
isolados -- como base para a geração de respostas.

\subsubsection{\textit{Pipeline} de indexação}
\label{sec-fundteo-graphrag-indexacao}

A indexação de um corpus no GraphRAG é organizada como um \textit{pipeline}
de \textit{workflows} executados sequencialmente, cada um responsável por uma
etapa específica do processamento:

\begin{itemize}
    \item \textbf{Carregamento e segmentação do texto:} os documentos de
    entrada são carregados e divididos em unidades de texto
    (\textit{text units}) de tamanho configurável, medido em número de
    \textit{tokens}, com sobreposição (\textit{overlap}) entre unidades
    consecutivas para preservar contexto nas bordas dos fragmentos.

    \item \textbf{Extração do grafo:} para cada unidade de texto, o LLM
    configurado é utilizado para extrair entidades e relações entre elas,
    segundo uma taxonomia de tipos de entidade fornecida na configuração
    (por exemplo, organização, pessoa, evento -- no caso deste trabalho,
    adaptada para o domínio jurídico). O processo suporta múltiplas
    ``passadas'' de extração sobre a mesma unidade de texto (parâmetro
    \textit{gleanings}), nas quais o modelo é solicitado a identificar
    entidades adicionais que não tenham sido capturadas na primeira
    passagem, aumentando a cobertura da extração.

    \item \textbf{Resumo de descrições:} quando a mesma entidade é
    mencionada em diferentes unidades de texto, suas descrições parciais são
    consolidadas em uma única descrição coerente, também com auxílio do LLM.

    \item \textbf{Detecção de comunidades:} o grafo de entidades e relações
    extraído é submetido a um algoritmo de detecção de comunidades -- o
    GraphRAG utiliza, por padrão, o algoritmo de Leiden \cite{traag2019louvain} --, que particiona o
    grafo em grupos hierárquicos de entidades fortemente conectadas entre si,
    formando diferentes níveis de granularidade (de comunidades menores e mais
    específicas a comunidades maiores e mais abrangentes).

    \item \textbf{Geração de relatórios de comunidade:} para cada comunidade
    identificada, o LLM gera um relatório textual que resume as entidades,
    relações e temas predominantes daquele agrupamento. Esses relatórios são,
    posteriormente, a base para a busca global.

    \item \textbf{Geração de \textit{embeddings}:} por fim, são gerados
    vetores de \textit{embedding} para as descrições de entidades, para as
    unidades de texto originais e para os relatórios de comunidade,
    armazenados em um banco de dados vetorial (por padrão, LanceDB \cite{lancedb}) para uso
    nas etapas de consulta que combinam busca vetorial com a estrutura do
    grafo.
\end{itemize}

Toda a saída do \textit{pipeline} de indexação -- entidades, relações,
comunidades, relatórios e unidades de texto -- é persistida em arquivos no
formato \textit{parquet}, o que facilita tanto o consumo por outras aplicações
(como a aplicação web deste trabalho) quanto a inspeção manual dos dados
durante o desenvolvimento.

\subsubsection{Modalidades de consulta}
\label{sec-fundteo-graphrag-consultas}

Uma das características centrais do GraphRAG, e de particular interesse para
este trabalho, é a disponibilização de diferentes modalidades de busca, cada
uma adequada a um tipo distinto de pergunta:

\begin{itemize}
    \item \textbf{Busca local (\textit{Local Search}):} indicada para
    perguntas pontuais, focadas em entidades específicas (por exemplo, ``o
    que diz o artigo 18 do CDC sobre vício do produto?''). Combina dados
    estruturados do grafo (entidades relacionadas, suas descrições e relações
    diretas) com trechos de texto não estruturado associados a essas
    entidades, compondo um contexto rico e específico para o LLM gerador.

    \item \textbf{Busca global (\textit{Global Search}):} indicada para
    perguntas amplas ou holísticas, que exigem sintetizar informação
    distribuída por diversas partes do corpus (por exemplo, ``quais são os
    principais direitos do consumidor previstos no CDC?''). Opera sobre os
    relatórios de comunidade gerados na indexação, utilizando uma estratégia
    de \textit{map-reduce}: cada relatório de comunidade é processado
    individualmente para gerar uma resposta parcial (\textit{map}), e essas
    respostas parciais são então combinadas em uma resposta final consolidada
    (\textit{reduce}).

    \item \textbf{Busca \textit{DRIFT} (\textit{Dynamic Reasoning and
    Inference with Flexible Traversal}):} combina elementos das buscas local
    e global de forma iterativa. Parte de uma resposta inicial abrangente
    (\textit{primer}), gerada de forma semelhante à busca global, e a refina
    por meio de sucessivas perguntas de acompanhamento (\textit{follow-ups}),
    respondidas via busca local, permitindo maior profundidade de raciocínio
    sobre a pergunta original.

    \item \textbf{Busca básica (\textit{Basic Search}):} realiza busca
    vetorial simples sobre as unidades de texto originais, sem envolver a
    estrutura do grafo ou os relatórios de comunidade. Funciona, na prática,
    como uma implementação de RAG vetorial tradicional dentro do próprio
    GraphRAG, compartilhando o mesmo corpus indexado, modelo de
    \textit{embeddings} e modelo gerador utilizados nas demais modalidades.
\end{itemize}

\subsubsection{Aspectos de configuração e execução}
\label{sec-fundteo-graphrag-config}

Em termos de implementação, a biblioteca GraphRAG é configurada por meio de um
arquivo \texttt{settings.yaml}, no qual são definidos os modelos de linguagem
utilizados (um para completions/geração e outro para \textit{embeddings}),
parâmetros de segmentação do texto, a taxonomia de tipos de entidade a ser
extraída, parâmetros de particionamento de comunidades e as configurações
específicas de cada modalidade de consulta. A comunicação com os provedores de
modelos de linguagem é abstraída pela biblioteca \textit{LiteLLM}, o que
permite configurar o GraphRAG para utilizar tanto provedores comerciais
baseados em API (como OpenAI, Google Gemini ou OpenRouter) quanto modelos
executados localmente por meio de ferramentas como o Ollama -- flexibilidade
explorada neste trabalho para viabilizar a execução do sistema sem
dependência obrigatória de custos de API. A biblioteca disponibiliza ainda um
utilitário de ajuste automático de \textit{prompts}
(\texttt{graphrag prompt-tune}), que gera exemplos de extração adaptados ao
vocabulário e à estrutura do corpus de entrada, reduzindo a necessidade de
engenharia manual de \textit{prompts} para domínios específicos como o
jurídico.

%%% Início de seção. %%%
\section{RAG Vetorial como Baseline de Comparação}
\label{sec-fundteo-baseline}

Um dos objetivos específicos deste trabalho consiste em avaliar
comparativamente a qualidade e a completude das respostas obtidas por meio do
GraphRAG frente a uma abordagem de RAG vetorial tradicional.Para isso, será utilizada 
como baseline a própria modalidade de \textbf{busca básica}
do GraphRAG (Seção~\ref{sec-fundteo-graphrag-consultas}).

Como essa modalidade realiza busca vetorial
pura sobre o mesmo corpus, com o mesmo modelo de \textit{embeddings} e o mesmo
modelo gerador utilizados nas demais modalidades, ela isola a variável de
interesse do experimento -- a presença ou ausência da estrutura de grafo --,
eliminando diferenças de configuração que poderiam confundir a comparação.


%%% Início de seção. %%%
\section{Tecnologias de Apoio ao Desenvolvimento da Aplicação Web}
\label{sec-fundteo-stack}

Como o foco principal deste trabalho é a investigação da abordagem GraphRAG
aplicada ao domínio jurídico, e não o desenvolvimento de uma aplicação web
sofisticada, a escolha das tecnologias de apoio priorizou simplicidade de
implementação e produtividade, de modo a permitir que a
maior parte do esforço de desenvolvimento seja dedicada à indexação, ao ajuste
do grafo de conhecimento e à avaliação dos resultados.

\subsection{\textit{Backend}}
\label{sec-fundteo-stack-backend}

O backend será implementado em Python utilizando o framework FastAPI,
devido à sua compatibilidade direta com a biblioteca GraphRAG, também
desenvolvida em Python. Essa integração permite invocar diretamente as
funcionalidades de consulta do framework, além de oferecer suporte a operações 
assíncronas e geração automática da documentação da API.

O \textit{backend} será responsável por expor endpoints REST para: (i)
submissão de consultas em linguagem natural, com seleção da modalidade de
busca desejada; (ii) obtenção dos dados do grafo de conhecimento (entidades,
relações e comunidades) para consumo pela interface de visualização.

\subsection{\textit{Frontend}}
\label{sec-fundteo-stack-frontend}

Para a interface web, propõe-se  um Single Page Application (SPA) utilizando React 
e TypeScript sobre o ecossistema de build Vite, garantindo agilidade e tipagem segura.
A estruturação visual e estilização adotaram o Tailwind CSS com componentes modulares 
baseados na biblioteca shadcn/ui, conferindo um design responsivo e com mínimo esforço de escrita de CSS.
O núcleo da aplicação — a visualização interativa do grafo gerado pelo GraphRAG — será construído com a biblioteca
react-force-graph-2d, que oferece uma exploração fluida, permitindo ao usuário navegar espacialmente (zoom e arraste),
alternar esquemas de cores por agrupamentos de comunidades, visualizar ferramentas de dicas (tooltips) em arestas e abrir
 painéis detalhados de entidades que exibem seus atributos e os relatórios contextuais das comunidades extraídas via LLM.

\subsection{Infraestrutura e Execução}
\label{sec-fundteo-stack-infra}

Todo o ambiente do projeto -- incluindo o \textit{backend}, o \textit{frontend},
o modelo de linguagem local (quando utilizado, por meio do Ollama) e o
\textit{pipeline} de indexação do GraphRAG -- será orquestrado por meio de
contêineres \textbf{Docker}, com definição dos serviços em um arquivo
\texttt{docker-compose.yml}. Essa escolha garante a reprodutibilidade do
ambiente de desenvolvimento e de execução, facilita a documentação do
processo de instalação no corpo do trabalho e permite que o sistema seja
executado de forma consistente independentemente do sistema operacional do
avaliador.